%
%
\section{Conclusion}




In this work, we propose State-Action Distillation (SAD), a novel approach for generating the pretraining dataset for ICRL, which is designed to overcome the limitations imposed by the existing ICRL algorithms like AD, DPT, and DIT in terms of relying on well-trained or even optimal policies to collect the pretraining dataset. 
%
SAD leverages solely random policies to construct the pretraining data, significantly promoting the practical application of ICRL in real-world scenarios. 
%
We also provide the quantitative analysis of the trustworthiness as well as the performance guarantees of SAD. Moreover, our empirical results on multiple popular ICRL benchmark environments demonstrate significant improvements over the existing baselines in terms of both performance and robustness.
%
Nevertheless, we note that SAD is currently limited to the discrete action space. Extending SAD to handle the continuous action space as well as more complex environments presents a promising direction for the future research.

%%%%%%%%%%%%%%%%%%%%%%%%%%%%%%%%%%%%%%%%%%%%%%%%%%%%%%%%%%%%%%%%%%%%%%%%

%%% The acknowledgments section is defined using the "acks" environment
%%% (rather than an unnumbered section). The use of this environment 
%%% ensures the proper identification of the section in the article 
%%% metadata as well as the consistent spelling of the heading.

% \begin{acks}
% If you wish to include any acknowledgments in your paper (e.g., to 
% people or funding agencies), please do so using the `\texttt{acks}' 
% environment. Note that the text of your acknowledgments will be omitted
% if you compile your document with the `\texttt{anonymous}' option.
% \end{acks}


